% claim.tex -- edgemid-monomer-dimer priority claim note (SOP 08b claim lane).
% AI-generated 2026-09-27 by the AI4Math paper department (claim lane).
% lstlean.tex copied unchanged from harness/templates/paper/lstlean.tex.
% Machine-readable anchors: every theorem carries a % LEAN: line (task id +
% declaration name + grade); paper-lint.sh and the claim audit check against it.
% The Lean listing is a byte-exact mechanical copy of the frozen final
% development tasks/20260926-edgemid-pipeline/final/03-proof-complete.lean
% (66 lines, sha256 1124e6a5...72c1d); its docstrings and header comment are
% in Chinese as frozen, rendered inside the listing via the package's
% CJK fontset (xeCJK/SimSun) plus an active-character rendering patch
% for the math symbols --- the sanctioned claims/cyl3-mod23 pattern
% (punctuation 00B7/2014/2026 routes through xeCJK, not the patch).
% Rendering-only; the listing body bytes stay identical to the frozen
% source (audit/listing-verify.txt).
\documentclass[11pt]{article}

\usepackage[T1]{fontenc}
\usepackage[utf8]{inputenc}
\usepackage{amsmath,amssymb,amsthm}
\usepackage{hyperref}
% lstlean.tex — Lean style for the listings package.
% Lineage: Jeremy Avigad (2015), adapted from Assia Mahboubi's SSR style;
% inspired by Olivier Verdier's unixode.sty for the Unicode literate table.
% No CTAN package or upstream repo exists; papers carry this file in their
% source package (cap set arXiv:1907.01449, sphere eversion arXiv:2210.07746).
% AI4Math adaptation (AI-generated, 2026-09-26): sorry/admit/sorryAx elevated
% to a dedicated error tier, matching the pipeline's 0-sorry constitution.
\usepackage{listings}
\usepackage{xcolor}

\definecolor{keywordcolor}{rgb}{0.7, 0.1, 0.1}   % red keywords
\definecolor{tacticcolor}{rgb}{0.0, 0.1, 0.6}    % blue tactics
\definecolor{commentcolor}{rgb}{0.4, 0.4, 0.4}   % grey comments
\definecolor{symbolcolor}{rgb}{0.0, 0.1, 0.6}    % blue symbols
\definecolor{sortcolor}{rgb}{0.1, 0.5, 0.1}      % green sorts
\definecolor{errorcolor}{rgb}{0.6, 0, 0}         % dark red: sorry/admit
\definecolor{stringcolor}{rgb}{0.5, 0.1, 0.5}    % purple strings

\lstdefinelanguage{lean}{
  % Anything between $ becomes LaTeX math mode
  mathescape=false,
  % Comments may or not include Latex commands
  texcl=false,
  % keywords, list taken from lean-syntax.txt
  morekeywords=[1]{import, prelude, protected, private, noncomputable,
    definition, meta, renaming, hiding, parameter, parameters, begin,
    constant, constants, lemma, variable, variables, theory, print,
    theorem, example, open, section, namespace, universe, universes, end,
    modifier, modifiers, using, export, exposing, axiom, inductive,
    coinductive, with, structure, class, instance, attribute, local, where,
    by, have, show, let, in, if, then, else, match, do, for, fun, assume,
    from, take, calc, include, omit, infix, infixl, infixr, notation,
    macro, syntax, elab, set_option, initialize, deriving, opaque, abbrev,
    mutual, partial, scoped, termination_by, decreasing_by},
  % Sorts
  morekeywords=[2]{Sort, Type, Prop},
  % tactics, list taken from lean-syntax.txt (tactic keywords)
  morekeywords=[3]{intro, intros, apply, exact, refine, rfl, simp, simp_all,
    simpa, simp_rw, rw, rewrite, erewrite, ring, ring_nf, omega, decide,
    norm_num, norm_cast, induction, cases, rcases, obtain, constructor,
    ext, funext, conv, congr, field_simp, linarith, nlinarith, positivity,
    tauto, trivial, aesop, use, by_cases, by_contra, push_neg, contrapose,
    symm, trans, assumption, contradiction, exfalso, specialize,
    generalize, revert, clear, rename_i, subst, unfold, delta, change,
    convert, split_ifs, interval_cases, fin_cases, zify, gcongr,
    nth_rewrite, suffices, generalizing, at, only, native_decide},
  % warnings - constitution tier: must never survive into a final proof
  morekeywords=[4]{sorry, admit, sorryAx},
  literate=
    {α}{{\ensuremath{\alpha}}}1
    {β}{{\ensuremath{\beta}}}1
    {γ}{{\ensuremath{\gamma}}}1
    {δ}{{\ensuremath{\delta}}}1
    {ε}{{\ensuremath{\varepsilon}}}1
    {ζ}{{\ensuremath{\zeta}}}1
    {η}{{\ensuremath{\eta}}}1
    {θ}{{\ensuremath{\theta}}}1
    {ι}{{\ensuremath{\iota}}}1
    {κ}{{\ensuremath{\kappa}}}1
    {λ}{{\ensuremath{\lambda}}}1
    {μ}{{\ensuremath{\mu}}}1
    {ν}{{\ensuremath{\nu}}}1
    {ξ}{{\ensuremath{\xi}}}1
    {π}{{\ensuremath{\pi}}}1
    {ρ}{{\ensuremath{\rho}}}1
    {σ}{{\ensuremath{\sigma}}}1
    {τ}{{\ensuremath{\tau}}}1
    {υ}{{\ensuremath{\upsilon}}}1
    {φ}{{\ensuremath{\varphi}}}1
    {χ}{{\ensuremath{\chi}}}1
    {ψ}{{\ensuremath{\psi}}}1
    {ω}{{\ensuremath{\omega}}}1
    {Γ}{{\ensuremath{\Gamma}}}1
    {Δ}{{\ensuremath{\Delta}}}1
    {Θ}{{\ensuremath{\Theta}}}1
    {Λ}{{\ensuremath{\Lambda}}}1
    {Ξ}{{\ensuremath{\Xi}}}1
    {Π}{{\ensuremath{\Pi}}}1
    {Σ}{{\ensuremath{\Sigma}}}1
    {Φ}{{\ensuremath{\Phi}}}1
    {Ψ}{{\ensuremath{\Psi}}}1
    {Ω}{{\ensuremath{\Omega}}}1
    {ℵ}{{\ensuremath{\aleph}}}1
    {∀}{{\ensuremath{\forall}}}1
    {∃!}{{\ensuremath{\exists}!}}1
    {∃}{{\ensuremath{\exists}}}1
    {∈}{{\ensuremath{\in}}}1
    {∉}{{\ensuremath{\notin}}}1
    {∘}{{\ensuremath{\circ}}}1
    {∩}{{\ensuremath{\cap}}}1
    {∪}{{\ensuremath{\cup}}}1
    {⊆}{{\ensuremath{\subseteq}}}1
    {⊂}{{\ensuremath{\subset}}}1
    {⊇}{{\ensuremath{\supseteq}}}1
    {⊃}{{\ensuremath{\supset}}}1
    {⊄}{{\ensuremath{\nsubseteq}}}1
    {∅}{{\ensuremath{\emptyset}}}1
    {∧}{{\ensuremath{\wedge}}}1
    {∨}{{\ensuremath{\vee}}}1
    {¬}{{\ensuremath{\neg}}}1
    {⊤}{{\ensuremath{\top}}}1
    {⊥}{{\ensuremath{\bot}}}1
    {↔}{{\ensuremath{\leftrightarrow}}}1
    {→}{{\ensuremath{\rightarrow}}}1
    {←}{{\ensuremath{\leftarrow}}}1
    {↑}{{\ensuremath{\uparrow}}}1
    {⇒}{{\ensuremath{\Rightarrow}}}1
    {⇐}{{\ensuremath{\Leftarrow}}}1
    {⟹}{{\ensuremath{\Longrightarrow}}}1
    {↦}{{\ensuremath{\mapsto}}}1
    {≤}{{\ensuremath{\leq}}}1
    {≥}{{\ensuremath{\geq}}}1
    {≠}{{\ensuremath{\neq}}}1
    {≈}{{\ensuremath{\approx}}}1
    {≡}{{\ensuremath{\equiv}}}1
    {≃}{{\ensuremath{\simeq}}}1
    {≅}{{\ensuremath{\cong}}}1
    {×}{{\ensuremath{\times}}}1
    {·}{{\ensuremath{\cdot}}}1
    {±}{{\ensuremath{\pm}}}1
    {⊕}{{\ensuremath{\oplus}}}1
    {⊗}{{\ensuremath{\otimes}}}1
    {⋆}{{\ensuremath{\star}}}1
    {▸}{{\ensuremath{\blacktriangleright}}}1
    {⟨}{{\ensuremath{\langle}}}1
    {⟩}{{\ensuremath{\rangle}}}1
    {⦃}{{\ensuremath{\{\!\mid}}}1
    {⦄}{{\ensuremath{\mid\!\}}}}1
    {⟦}{{\ensuremath{[\![}}}1
    {⟧}{{\ensuremath{]\!]}}}1
    {⌊}{{\ensuremath{\lfloor}}}1
    {⌋}{{\ensuremath{\rfloor}}}1
    {⌈}{{\ensuremath{\lceil}}}1
    {⌉}{{\ensuremath{\rceil}}}1
    {√}{{\ensuremath{\surd}}}1
    {∣}{{\ensuremath{\mid}}}1
    {∤}{{\ensuremath{\nmid}}}1
    {∎}{{\ensuremath{\blacksquare}}}1
    {⋯}{{\ensuremath{\cdots}}}1
    {…}{{\ensuremath{\ldots}}}1
    {∑}{{\ensuremath{\sum}}}1
    {∏}{{\ensuremath{\prod}}}1
    {⁻¹}{{\ensuremath{^{-1}}}}1
    {¹}{{\ensuremath{^1}}}1
    {²}{{\ensuremath{^2}}}1
    {³}{{\ensuremath{^3}}}1
    {ⁿ}{{\ensuremath{^n}}}1
    {ᵐ}{{\ensuremath{^m}}}1
    {₀}{{\ensuremath{_0}}}1
    {₁}{{\ensuremath{_1}}}1
    {₂}{{\ensuremath{_2}}}1
    {₃}{{\ensuremath{_3}}}1
    {₄}{{\ensuremath{_4}}}1
    {₅}{{\ensuremath{_5}}}1
    {ₙ}{{\ensuremath{_n}}}1
    {ₘ}{{\ensuremath{_m}}}1
    {ₖ}{{\ensuremath{_k}}}1
    {ᵢ}{{\ensuremath{_i}}}1
    {ⱼ}{{\ensuremath{_j}}}1
    {ℤ}{{\ensuremath{\mathbb{Z}}}}1
    {ℕ}{{\ensuremath{\mathbb{N}}}}1
    {ℚ}{{\ensuremath{\mathbb{Q}}}}1
    {ℝ}{{\ensuremath{\mathbb{R}}}}1
    {ℂ}{{\ensuremath{\mathbb{C}}}}1
    {𝔽}{{\ensuremath{\mathbb{F}}}}1
    {𝒫}{{\ensuremath{\mathcal{P}}}}1
    {∗}{{\ensuremath{\ast}}}1
    {•}{{\ensuremath{\bullet}}}1,
  % Comments
  morecomment=[l]{--},
  morecomment=[s]{/-}{-/},
  morestring=[b]",
  stringstyle=\color{stringcolor},
  keepspaces=true,
  keywordstyle=[1]\color{keywordcolor}\bfseries,
  keywordstyle=[2]\color{sortcolor}\bfseries,
  keywordstyle=[3]\color{tacticcolor}\bfseries,
  keywordstyle=[4]\color{errorcolor}\bfseries\underline,
  escapeinside={(*@}{@*)},
  columns=fullflexible,
  basicstyle=\ttfamily\small,
  commentstyle=\color{commentcolor}\itshape,
  breaklines=true,
  showstringspaces=false,
  frame=single,
  framesep=2mm,
  xleftmargin=4mm,
  numbers=left,
  numberstyle=\tiny\color{commentcolor},
  numbersep=6pt,
}

% Inline Lean code in prose: \lean{Int.ModEq.pow}
\newcommand{\lean}[1]{\lstinline[language=lean,basicstyle=\ttfamily\normalsize]|#1|}

\lstset{language=lean}
% --- XeTeX rendering patch (verbatim from claims/tatami-mod8-defect/claim.tex
%     2026-09-27 sanctioned block, itself backported from
%     papers/cyl3-mod23/paper.tex; extended with the 6 further math symbols
%     of this frozen file's inventory: 00D7 2083 2099 2212 2223 27FA;
%     00B7/2014/2026 are in xeCJK's punctuation lists (MiddlePunct/LongPunct,
%     xeCJK.sty) and are left to xeCJK's own CJK font routing).
%     Rendering-only; source bytes of listings unchanged. Guarded so pdflatex
%     keeps literate-table behaviour.
\ifdefined\XeTeXversion
  {\lccode`\~="2022 \lowercase{\global\catcode`~=\active \global\def~{\ensuremath{\bullet}}}}
  {\lccode`\~="2115 \lowercase{\global\catcode`~=\active \global\def~{\ensuremath{\mathbb{N}}}}}
  {\lccode`\~="2124 \lowercase{\global\catcode`~=\active \global\def~{\ensuremath{\mathbb{Z}}}}}
  {\lccode`\~="2190 \lowercase{\global\catcode`~=\active \global\def~{\ensuremath{\leftarrow}}}}
  {\lccode`\~="2192 \lowercase{\global\catcode`~=\active \global\def~{\ensuremath{\rightarrow}}}}
  {\lccode`\~="2194 \lowercase{\global\catcode`~=\active \global\def~{\ensuremath{\leftrightarrow}}}}
  {\lccode`\~="2200 \lowercase{\global\catcode`~=\active \global\def~{\ensuremath{\forall}}}}
  {\lccode`\~="2203 \lowercase{\global\catcode`~=\active \global\def~{\ensuremath{\exists}}}}
  {\lccode`\~="2227 \lowercase{\global\catcode`~=\active \global\def~{\ensuremath{\wedge}}}}
  {\lccode`\~="2228 \lowercase{\global\catcode`~=\active \global\def~{\ensuremath{\vee}}}}
  {\lccode`\~="2261 \lowercase{\global\catcode`~=\active \global\def~{\ensuremath{\equiv}}}}
  {\lccode`\~="2264 \lowercase{\global\catcode`~=\active \global\def~{\ensuremath{\leq}}}}
  {\lccode`\~="2265 \lowercase{\global\catcode`~=\active \global\def~{\ensuremath{\geq}}}}
  {\lccode`\~="22A2 \lowercase{\global\catcode`~=\active \global\def~{\ensuremath{\vdash}}}}
  {\lccode`\~="27E8 \lowercase{\global\catcode`~=\active \global\def~{\ensuremath{\langle}}}}
  {\lccode`\~="27E9 \lowercase{\global\catcode`~=\active \global\def~{\ensuremath{\rangle}}}}
  {\lccode`\~="00D7 \lowercase{\global\catcode`~=\active \global\def~{\ensuremath{\times}}}}
  {\lccode`\~="2083 \lowercase{\global\catcode`~=\active \global\def~{\ensuremath{{}_3}}}}
  {\lccode`\~="2099 \lowercase{\global\catcode`~=\active \global\def~{\ensuremath{{}_n}}}}
  {\lccode`\~="2212 \lowercase{\global\catcode`~=\active \global\def~{\ensuremath{-}}}}
  {\lccode`\~="2223 \lowercase{\global\catcode`~=\active \global\def~{\ensuremath{\mid}}}}
  {\lccode`\~="27FA \lowercase{\global\catcode`~=\active \global\def~{\ensuremath{\Longleftrightarrow}}}}
\fi
\ifdefined\XeTeXversion
  \usepackage{xeCJK}
  \setCJKmainfont{SimSun}
\fi

\newtheorem{theorem}{Theorem}

\title{Mod-2 periodicity for a family of sixth-order linear recurrences,
with a matrix-functional generalization\\
(Lean 4 machine-checked claim of record)}
\author{Aurora0134\thanks{Contact: via the GitHub account Aurora0134; ORCID to
be added at Zenodo upload.}}
\date{September 27, 2026}

\begin{document}
\maketitle

\begin{abstract}
We record two theorems about the sixth-order linear recurrence
$c(n+6)=4c(n+5)+14c(n+4)-10c(n+2)+c(n)$. First, every integer sequence
satisfying it has parity of period dividing $6$, and the parity of $c(n)$ is
determined by $n \bmod 6$. Second, a matrix-functional generalization: any
$8\times 8$ integer matrix $T$ annihilated by $x^6-4x^5-14x^4+10x^2-1$ yields
scalar sequences $n\mapsto (T^n v)_i$ obeying the same recurrence. The
motivation comes from monomer--dimer (matching) enumeration on $3\times n$
boards with one boundary defect, where the full-board counts are OEIS A033506
and the defect-family sequences behind this recurrence are recorded here for
the first time; that combinatorial interpretation is \emph{not} part of the
machine-checked claim. All proofs are machine-checked by the Lean 4 kernel
(0 \texttt{sorry}; axioms within \{propext, Quot.sound, Classical.choice\}). A
structured novelty search found no occupying record. The artifact package is
deposited on Zenodo. This note is a priority claim of record, not a full
paper.
\end{abstract}

\section{Introduction}

Counting monomer--dimer coverings of a $3\times n$ board --- matchings of the
grid graph $P_3 \times P_n$ --- is a classical enumeration problem: the
full-board counts form OEIS sequence A033506 \cite{oeis-a033506}, with a
generating function documented there following Lundow's enumeration of
matchings in polygraphs \cite{lundow-1996,lundow-1998}; see Propp's survey
for the landscape \cite{propp-1999}. Forcing one prescribed boundary cell to
carry a monomer (deleting one vertex, in matching language) produces
\emph{defect} counting sequences. The pipeline's transfer-matrix computations
of three such families --- a side-middle hole, a top-edge hole, and a fixed
edge cut --- produce integer sequences that, as of September 2026, do not
appear in the OEIS under calibrated exact-term searches, and whose recurrences
share the sixth-order signature $(4,14,0,-10,0,1)$ of the full board. The
nearest literature studies \emph{perfect} matchings of vertex-deleted lattices
(pure dimer coverings), on odd-by-odd boards \cite{tzeng-wu-2003,wu-2006,oh-2019};
adjacent non-occupying works prove a reciprocity theorem for total matchings
of full boards \cite{anzalone-2003} and study the matching complex of the same
$3\times n$ family topologically \cite{goyal-2021}.

This note records two theorems at the recurrence level, produced by AI4Math, a
highly automated LLM-plus-kernel pipeline in which the Lean 4 kernel is the
sole arbiter of correctness. The present text is a priority claim of record
--- a short deposit note accompanying a Zenodo artifact package; a fuller
narrative paper is prepared separately through the pipeline's paper lane.
Every theorem stated below is a \emph{complete proof} in the pipeline's
grading: it compiles with 0 \texttt{sorry}, its axioms stay within the
standard whitelist, and the statement was frozen before proving. The
combinatorial reading of the recurrence (tiling counts) is motivation only;
Section~\ref{sec:scope} is explicit about what the kernel did \emph{not}
prove.

\section{Statement}

The formal development adds no definitions: each theorem is self-contained,
quantifying over the objects in its hypothesis. The frozen statement files
contain exactly one theorem each (\texttt{edgemid\_master} in
\texttt{edgemid-p1-statement.lean}, sha256 \texttt{f964ce71...27b27c56};
\texttt{edgemid\_matrix\_family} in \texttt{edgemid-p2-statement.lean}, sha256
\texttt{cf9e379a...99b6d5}). The formal theorem headers, including their
Chinese docstrings as frozen, appear verbatim inside Listing~\ref{lst:full}
(lines 15--20 and 44--52 there), byte-identical to the frozen files.

% LEAN: tasks/20260926-edgemid-pipeline :: edgemid_master grade=完全证明
\begin{theorem}[mod-2 periodicity and characterization]\label{thm:p1}
Let $c:\mathbb{N}\to\mathbb{Z}$ satisfy
$c(n+6)=4c(n+5)+14c(n+4)-10c(n+2)+c(n)$ for all $n\in\mathbb{N}$. Then:
\begin{enumerate}
\item $c(n+6)\equiv c(n)\pmod 2$ for all $n$; and
\item $c(n)\equiv c(n\bmod 6)\pmod 2$ for all $n$, i.e.\ the parity of $c(n)$
is completely determined by $n \bmod 6$.
\end{enumerate}
\end{theorem}

% LEAN: tasks/20260926-edgemid-pipeline :: edgemid_matrix_family grade=完全证明
\begin{theorem}[matrix-functional generalization]\label{thm:p2}
Let $T$ be an $8\times 8$ integer matrix satisfying
$T^6=4T^5+14T^4-10T^2+I$, let $v\in\mathbb{Z}^8$ and let $i$ be an index.
Then for all $n\in\mathbb{N}$,
\[
(T^{n+6}v)_i \;=\; 4\,(T^{n+5}v)_i + 14\,(T^{n+4}v)_i - 10\,(T^{n+2}v)_i + (T^n v)_i,
\]
i.e.\ the scalar sequence $n\mapsto (T^n v)_i$ satisfies the recurrence of
Theorem~\ref{thm:p1}.
\end{theorem}

\section{Proof}

The complete machine proofs are reproduced verbatim in
Listing~\ref{lst:full}, which is the whole frozen final development
\texttt{final/03-proof-complete.lean} (66 lines, sha256
\texttt{1124e6a5...72c1d}): imports, header comment, both theorems with their
full proofs. In brief, Theorem~\ref{thm:p1} converts the congruence into a
divisibility via \lean{Int.modEq_iff_dvd}, discharges the parity
rearrangement with the recurrence hypothesis and \texttt{ring}, lifts the
one-step congruence to six-step blocks by induction
(\lean{Int.ModEq.refl}, \lean{Int.ModEq.trans}), and finishes with the
division identity $n = n\bmod 6 + 6\lfloor n/6\rfloor$ (\texttt{omega}).
Theorem~\ref{thm:p2} rewrites $T^{n+6}=T^n T^6$ (\texttt{pow\_add}),
substitutes the annihilating relation, distributes over the matrix operations
(\lean{Matrix.add_mulVec}, \lean{Matrix.sub_mulVec},
\lean{Matrix.smul_mulVec}, \lean{mul_smul_comm}), and closes by \texttt{rfl}
along a finite fallback chain. A narrative prose proof is not the duty of
this note; it belongs to the full paper lane.

\begin{lstlisting}[caption={Complete frozen final development (verbatim full text, lines 1--66 of \texttt{final/03-proof-complete.lean}, sha256 \texttt{1124e6a5...72c1d}). The two theorem headers (lines 15--20, 44--52) are byte-identical to the frozen statement files; docstrings and the header comment are in Chinese as frozen.},label={lst:full}]
import Mathlib.Data.Int.ModEq
import Mathlib.Data.Matrix.Mul

/-!
# 03 证明完成稿（候选 03 · edgemid-pipeline · P1+P2）

两定理 statement 头与冻结快照逐字一致（宪条 2）：P1 锚 `formalized/edgemid-p1-statement.lean:12-17`，
P2 锚 `formalized/edgemid-p2-statement.lean:14-22`（两者均锚定
`tasks/20260926-beian-extselect/candidates/smoke-03/EdgemidSmoke03r3.lean` 定理头）。
proof body 逐字取自 smoke r3（P1 :18-39 / P2 :50-63，逐字节一致零改动），
全稿零占位符、无 namespace。imports = 两冻结件 import 面并集。
AI 生成（2026-09-26 军政部工人，dynamic-workflow 子代理）。
-/

/-- P1(首选命题):凡满足六阶线性递推
`c (n+6) = 4 * c (n+5) + 14 * c (n+4) - 10 * c (n+2) + c n` 的整数序列,
模 2 以 6 为周期,且模 2 余数由 `n % 6` 完全决定(整除刻画支)。 -/
theorem edgemid_master (c : ℕ → ℤ)
    (h : ∀ n, c (n+6) = 4 * c (n+5) + 14 * c (n+4) - 10 * c (n+2) + c n) :
    (∀ n, c (n+6) ≡ c n [ZMOD 2]) ∧ (∀ n, c n ≡ c (n % 6) [ZMOD 2]) := by
  have key : ∀ n, c (n+6) ≡ c n [ZMOD 2] := by
    intro n
    refine Int.modEq_iff_dvd.2 ⟨-(2 * c (n+5) + 7 * c (n+4) - 5 * c (n+2)), ?_⟩
    rw [h n]
    ring
  have aux : ∀ k n, c (n + 6 * k) ≡ c n [ZMOD 2] := by
    intro k
    induction k with
    | zero =>
        intro n
        simp only [Nat.mul_zero, Nat.add_zero]
        exact Int.ModEq.refl _
    | succ k ih =>
        intro n
        have e : n + 6 * (k+1) = (n + 6 * k) + 6 := by ring
        rw [e]
        exact (key (n + 6 * k)).trans (ih n)
  refine ⟨key, fun n => ?_⟩
  have step : c (n % 6 + 6 * (n / 6)) ≡ c (n % 6) [ZMOD 2] := aux (n / 6) (n % 6)
  have e : n % 6 + 6 * (n / 6) = n := by omega
  rw [e] at step
  exact step

/-- P2(挑战档扩展):转移矩阵被六阶多项式零化时,其矩阵幂泛函
(含各固定孔位缺陷计数)满足同一递推。 -/
theorem edgemid_matrix_family {T : Matrix (Fin 8) (Fin 8) ℤ}
    (v : Fin 8 → ℤ) (i : Fin 8)
    (hT : T ^ 6 = 4 • T ^ 5 + 14 • T ^ 4 - 10 • T ^ 2 + (1 : Matrix (Fin 8) (Fin 8) ℤ))
    (n : ℕ) :
    (T ^ (n+6)).mulVec v i
      = 4 * (T ^ (n+5)).mulVec v i + 14 * (T ^ (n+4)).mulVec v i
        - 10 * (T ^ (n+2)).mulVec v i + (T ^ n).mulVec v i := by
  have hgen : T ^ (n+6)
      = 4 • T ^ (n+5) + 14 • T ^ (n+4) - 10 • T ^ (n+2) + (T^n : Matrix (Fin 8) (Fin 8) ℤ) := by
    rw [pow_add, hT, Matrix.mul_add, Matrix.mul_sub, Matrix.mul_add,
        mul_smul_comm, mul_smul_comm, mul_smul_comm,
        Matrix.mul_one, ← pow_add, ← pow_add, ← pow_add]
  rw [hgen, Matrix.add_mulVec, Matrix.sub_mulVec, Matrix.add_mulVec,
      Matrix.smul_mulVec, Matrix.smul_mulVec, Matrix.smul_mulVec]
  first
  | rfl
  | simp only [smul_eq_mul]
  | simp [smul_eq_mul]
  | norm_num
  | push_cast <;> rfl
  | lia
\end{lstlisting}

\section{Verification evidence}

\begin{itemize}
\item Toolchain: Lean v4.34.0 (\texttt{leanprover/lean4:v4.34.0})
\cite{lean4}, mathlib4 v4.34.0 (rev \texttt{5ed29652}) \cite{mathlib}.
\item Statement integrity: frozen statements
\texttt{formalized/edgemid-p1-statement.lean}, sha256
\texttt{f964ce71e9888d7e8d9d982996fed2531492afdb1ca47df004a7f5d727b27c56},
and \texttt{formalized/edgemid-p2-statement.lean}, sha256
\texttt{cf9e379af46affeee0306ba2bb84094a18507eaeacbae9d0b690ac71cb99b6d5};
final development \texttt{final/03-proof-complete.lean}, sha256
\texttt{1124e6a53b6ec70fee5c964490fa4f407a23ab03b1557d3f7d214b0c20572c1d};
symmetric difference of the two theorem headers against the frozen snapshots
(frozen lines 12--17 vs final lines 15--20; frozen lines 14--22 vs final
lines 44--52): zero, byte-identical.
\item Kernel check: strict compile exit 0 with no errors; a byte-level scan of
the whole final file (comments included) finds 0 occurrences of
\texttt{sorry}/\texttt{admit}. Gate-2 non-degeneracy: all seven bare-tactic
probes (\texttt{nlinarith}, \texttt{norm\_num}, \texttt{simp},
\texttt{ring\_nf}, \texttt{positivity}, \texttt{omega}, \texttt{linarith})
fail to close either goal. Axiom audit, verbatim from
\texttt{audit/final-axioms-p1.txt} and \texttt{audit/final-axioms-p2.txt}:
\begin{verbatim}
WELLDEF OK
AXIOMS edgemid_master [propext, Quot.sound]
\end{verbatim}
\begin{verbatim}
WELLDEF OK
AXIOMS edgemid_matrix_family [propext, Classical.choice, Quot.sound]
\end{verbatim}
Both lists are subsets of the whitelist \{propext, Quot.sound,
Classical.choice\}; \texttt{sorryAx} is absent.
\item Audit chain: statement freeze with sha256 snapshots; gate-2
well-definedness compilation plus degeneracy probing; proof assembly
restricted to proof bodies; gate-3 final strict compilation, byte-level
symmetric-difference comparison and axiom audit; an NL--Lean roundtrip
faithfulness walkthrough at formalization time. Evidence originals are in
the artifact package \texttt{audit/}.
\end{itemize}

\section{Novelty evidence}

The pipeline's C9 novelty review adjudicated this result as \emph{no
occupying record found}, and an independent exclusivity re-review
(2026-09-26, before deposit) maintained that adjudication. This is a
statement about search evidence, not a claim to new mathematics: the
full-board object A033506 is an occupied OEIS record \cite{oeis-a033506},
and the three defect-family sequences are recorded here for the first time at
the level of calibrated exact-term search evidence (wording capped
accordingly). Query channels and query-string summary (the full query log,
including zero-hit evidence files, is in \texttt{audit/c9-record.md} of the
package, copied verbatim from the source record):

\begin{center}
\begin{tabular}{p{0.20\textwidth}p{0.50\textwidth}p{0.21\textwidth}}
\hline
channel & query summary & result \\
\hline
OEIS sequence layer & calibration: baseline 8 terms
$1,3,22,131,823,5096,31687,196785$ & hits A033506 (channel works) \\
OEIS sequence layer & exact 8-term W / N1 / EB0; relaxed mid-window 5-term
variants; corner-hole control 8 terms & all null (zero hits) \\
OEIS record layer & six full-text queries (monomer-dimer; deleted vertex +
matchings; grid graph + matchings; square removed + 3 x n; monomer + 3 x n;
matchings + one vertex); A033506 cross-references & hits are whole-board
families or tatami entries; no $3\times n$ defect variant \\
literature layer & arXiv (7 query groups, 8 abstracts verified, 1 full text
checked), OpenAlex (6 groups), zbMATH (5 groups), Semantic Scholar (1 success,
2 rate-limited) & no occupying record; nearest line
\cite{oh-2019,tzeng-wu-2003,wu-2006} demarcated, not occupying \\
web layer & general web, English and Chinese sides & whole-board classics and
OI/programming tutorials only \\
library layer & mathlib ripgrep for \texttt{monomer} / \texttt{dimer} & zero
hits \\
\hline
\end{tabular}
\end{center}

The nearest literature line --- Oh's state matrix recursion method
\cite{oh-2019} and the Tzeng--Wu / Wu Pfaffian results
\cite{tzeng-wu-2003,wu-2006} --- counts \emph{perfect} matchings of
vertex-deleted lattices (pure dimer coverings), the latter two restricted to
odd-by-odd boards; the sequences behind this note count \emph{all} matchings
of the vertex-deleted $3\times n$ graph, for all $n$. Different object, no
occupation; cited here for demarcation. Recorded residual gaps, verbatim in
spirit from the re-review: (i) CNKI (Chinese-language database) was not
searched (login required); (ii) two Semantic Scholar queries were
rate-limited (HTTP 429), with one query succeeding; (iii) the Tzeng--Wu 2003
and Wu 2006 original PDFs were not fetched --- the demarcation rests on the
ar5iv full text of \cite{oh-2019} plus abstracts.

\section{Scope and limitations}\label{sec:scope}

Grading boundary, per the final audit: the two theorems are complete proofs,
strictly as their Lean statements --- Theorem~\ref{thm:p1} about
\emph{arbitrary} integer sequences satisfying the recurrence, and
Theorem~\ref{thm:p2} as a uniform recurrence for matrix functionals of
\emph{arbitrary} annihilated $8\times 8$ integer matrices. Not
machine-checked here: (i) any identification of a concrete $T$ with the
transfer matrix of a $3\times n$ defect board, and any equality of
$(T^n v)_i$ with a defect matching count --- that bridge is the standard
transfer-matrix narrative \cite{oeis-a033506,lundow-1998,propp-1999}, and
mathlib currently ships no monomer--dimer or matching-enumeration API to
build it on (an eight-item gap list, ripgrep-verified, is in the artifact's
evidence pack); (ii) the claim that the W / N1 / EB0 defect-family sequences
satisfy the recurrence at all, which stands at the level of numerical
evidence (term-window fits) plus, for the full board, the OEIS-documented
generating function --- OEIS/literature endorsement and a computational
conjecture layer, not kernel content: the kernel proved no counting ontology.
The periodicity branch of Theorem~\ref{thm:p1} is algebraically shallow (a
one-line parity rearrangement); the characterization branch carries the
content --- stated to prevent over-reading. Novelty claims are search
evidence on a specific date, not proofs of absence. This note contains no
literature survey; it is a deposit record, not a survey article.

\section*{Data availability}

Artifact package: Zenodo deposit (DOI to be reserved at upload; see UPLOAD.md
in this package) with a public mirror on GitHub at
\url{https://github.com/Aurora0134/ai4math-claims}. Note under CC BY 4.0,
code under MIT.

\section*{Use of generative AI}

(a) Stage mapping: candidate-proposition generation, formalization into Lean
statements, retrieval evidence packing, proof assembly, audit documentation,
and text drafting (including this note) were performed by AI (LLM) stages of
the AI4Math pipeline; topic adjudication, statement-semantics verification
and gate sign-off were human. (b) Model/node/budget (numbers copied verbatim
from the source task ledger \texttt{tasks/20260926-edgemid-pipeline/budget.log}):
a single user-designated model node, wire-id
\texttt{anthropic/S3AI-Grok/grok-4.7} (model-detect resolution level L2;
session runtime GLM-5.3 Flash via official direct connection, per the
ledger's node-snapshot line), no cross-node switching. Source task totals:
llm-call $0/6$ (cap 6, expected 0 achieved --- all drafting ran on the
session runtime model through workflow subagents, so the call-metered budget
records zero and there are no pass-at-k rates to report); Lean compilations
$9/24$ on the kernel-script caliber ($18/24$ including worker-side
re-verification runs; both calibers are recorded line by line in the ledger);
diagnostics $0/4$; timebox about 2 hours 20 minutes within 24 hours. This
claim-note stage used llm-call $0$ of its budget $4$, and $7$ LaTeX
compilations of its budget $6$ (counting the final rebuild that embedded this
number; ledger \texttt{claims/edgemid-monomer-dimer/budget.log}). (c) Final
arbiter: all statements and proofs reported here were checked by the Lean 4
kernel: every proof compiles with no \texttt{sorry}, and \texttt{\#print
axioms} reports only subsets of \{propext, Quot.sound, Classical.choice\}.
(d) The human author reviewed and is responsible for all content; the AI
system is not an author.

\section*{Author contributions}

CRediT-style: the human author was responsible for topic selection and
adjudication, verification of statement semantics, gate sign-off, and the
decision to deposit; the AI4Math pipeline (an LLM system) performed
proposition generation, proof search, and text drafting. The AI system is
not listed as an author.

\begin{thebibliography}{9}
% Every entry below was verified on 2026-09-26: arXiv entries via the arXiv API
% (authors/title/date), journal entries via Crossref (DOI), OEIS and Lundow via
% the fetched OEIS entry text. No entry is from memory. Entries reused verbatim
% from the paper-lane manuscript papers/edgemid-monomer-dimer/paper.tex.
\bibitem{oeis-a033506}
The OEIS Foundation Inc.
\newblock The On-Line Encyclopedia of Integer Sequences, entry A033506:
\emph{Number of matchings in graph $P_3 \times P_n$}.
\newblock \url{https://oeis.org/A033506} (retrieved 2026-09-26).

\bibitem{lundow-1996}
P.~H. Lundow.
\newblock Computation of matching polynomials and the number of 1-factors in
polygraphs.
\newblock Research reports, No.~12, Department of Mathematics, Ume\r{a}
University, 1996. (As recorded in~\cite{oeis-a033506}.)

\bibitem{lundow-1998}
P.~H. Lundow.
\newblock Enumeration of matchings in polygraphs, 1998.
\newblock (As recorded in~\cite{oeis-a033506}.)

\bibitem{propp-1999}
J. Propp.
\newblock Enumeration of matchings: problems and progress.
\newblock In \emph{New Perspectives in Algebraic Combinatorics},
Math.\ Sci.\ Res.\ Inst.\ Publ.\ 38, Cambridge University Press, 1999.
\newblock \url{https://doi.org/10.1017/9781009701587.008}

\bibitem{tzeng-wu-2003}
W.-T. Tzeng and F.~Y. Wu.
\newblock Dimers on a simple-quartic net with a vacancy.
\newblock \emph{Journal of Statistical Physics} 110 (2003), 671--689.
\newblock \url{https://doi.org/10.1023/A:1022155701655}

\bibitem{wu-2006}
F.~Y. Wu.
\newblock Pfaffian solution of a dimer-monomer problem: single monomer on the
boundary.
\newblock \emph{Physical Review E} 74 (2006), 020104(R); erratum
\emph{ibid.} 74 (2006), 039907.
\newblock \url{https://doi.org/10.1103/PhysRevE.74.020104}

\bibitem{oh-2019}
S. Oh.
\newblock State matrix recursion method and monomer--dimer problem.
\newblock arXiv:1901.07847, 2019. \url{https://arxiv.org/abs/1901.07847}

\bibitem{anzalone-2003}
N. Anzalone, J. Baldwin, I. Bronshtein and T.~K. Petersen.
\newblock A reciprocity theorem for monomer-dimer coverings.
\newblock arXiv:math/0304359, 2003. \url{https://arxiv.org/abs/math/0304359}

\bibitem{goyal-2021}
S. Goyal, S. Shukla and A. Singh.
\newblock Matching complexes of $3 \times n$ grid graphs.
\newblock arXiv:2106.09915, 2021. \url{https://arxiv.org/abs/2106.09915}

\bibitem{lean4}
L.~de Moura and S.~Ullrich.
\newblock The Lean 4 theorem prover and programming language.
\newblock \emph{CADE-28}, 2021. \url{https://doi.org/10.1007/978-3-030-79876-5_27}

\bibitem{mathlib}
The mathlib Community.
\newblock The Lean mathematical library.
\newblock \emph{CPP 2020}, 2020. \url{https://doi.org/10.1145/3372885.3373824}
\end{thebibliography}

\end{document}
